\SourceLeaf{260}%
Substituting these measures of curvature into the series for \(\sigma\) gives the following expression, accurate through orders below the sixth:
\[
\begin{aligned}
\sigma=\tfrac{1}{2}ac\sin B\{&1+\tfrac{1}{120}\alpha(4pp-2qq+3qq'+3q'q')\\
&+\tfrac{1}{120}\beta(3pp-6qq+6qq'+3q'q')\\
&+\tfrac{1}{120}\gamma(3pp-2qq+qq'+4q'q')\}
\end{aligned}
\]
The precision is unchanged if we substitute for \(p,q,q'\) the quantities \(c\sin B,c\cos B,c\cos B-a\). We then obtain
\[
\begin{aligned}
\relax[8]\quad\sigma=\tfrac{1}{2}ac\sin B\{&1+\tfrac{1}{120}\alpha(3aa+4cc-9ac\cos B)\\
&+\tfrac{1}{120}\beta(3aa+3cc-12ac\cos B)\\
&+\tfrac{1}{120}\gamma(4aa+3cc-9ac\cos B)\}
\end{aligned}
\]
Since every term referring to the line \(AD\) drawn normal to \(BC\) has disappeared from this equation, we may also permute the points \(A,B,C\) and their corresponding quantities. Thus, with the same precision, we have
\[
\begin{aligned}
\relax[9]\quad\sigma=\tfrac{1}{2}bc\sin A\{&1+\tfrac{1}{120}\alpha(3bb+3cc-12bc\cos A)\\
&+\tfrac{1}{120}\beta(3bb+4cc-9bc\cos A)\\
&+\tfrac{1}{120}\gamma(4bb+3cc-9bc\cos A)\}
\end{aligned}
\]
\[
\begin{aligned}
\relax[10]\quad\sigma=\tfrac{1}{2}ab\sin C\{&1+\tfrac{1}{120}\alpha(3aa+4bb-9ab\cos C)\\
&+\tfrac{1}{120}\beta(4aa+3bb-9ab\cos C)\\
&+\tfrac{1}{120}\gamma(3aa+3bb-12ab\cos C)\}
\end{aligned}
\]
\Article{26}
It is very useful to consider a plane rectilinear triangle with sides equal to \(a,b,c\). We shall denote its angles by \(A^*,B^*,C^*\). They differ from the angles of the triangle on the curved surface, namely \(A,B,C\), by quantities of the second order, and it is worth expanding these differences precisely. The calculations are more tedious than difficult, so it is enough to set down their main steps.

In formulas [1], [4], [5], replacing the quantities associated with \(B\) by their counterparts for \(C\) gives us formulas for \(r'r',r'\cos\varphi',r'\sin\varphi'\). Then expanding \(rr+r'r'-(q-q')^2-2r\cos\varphi\mathord{.}r'\cos\varphi'-2r\sin\varphi\mathord{.}r'\sin\varphi'\), which takes the form \(=bb+cc-aa-2bc\cos A=2bc(\cos A^*-\cos A)\), and combining that expansion with the expansion
